\section{The universal lower bound}\label{sec:optimality}

The examples constructed above attain the least possible cardinality and
weight. We prove this by an estimate for arbitrary abelian coefficient
sheaves, independent of the quotient construction. The relevant invariant
is the number of compact opens needed to cover a space.

Recall that \(P(W)\) denotes the integral free sheaf represented by an open
subset \(W\subseteq X\), so that

\[
\Hom(P(W),\mathcal M)\cong\Gamma(W,\mathcal M).
\]

In particular, \(P(X)\cong\underline{\Z}_X\). Throughout this section,
projective dimension is understood in its Ext-vanishing sense:
\(\pd(P)\leq n\) means that \(\Ext^d(P,M)=0\) for every coefficient object
\(M\) and every \(d>n\). The category of abelian sheaves has enough
injectives, so dimension shifting in the coefficient variable is available.
The indexing begins with countable covers: they already give projective
free-open sheaves, rather than merely a dimension bound of one. Each passage
to the next infinite cardinal costs at most one further degree. Keeping
this starting point explicit will give precisely the lower bound
\(\aleph_r\) in degree \(r\).

\subsection{Continuous transfinite extensions}

The homological input is a closure property whose proof is worth recording.
Continuity of the filtration is essential to its use below.

\begin{lemma}\label{lem:transfinite-dimension}
Let \(\mathcal A\) be an abelian category with enough injectives. Suppose
\((P_\alpha)_{\alpha\leq\lambda}\) is a well-ordered filtration with
\(P_0=0\), monomorphic successor maps, and

\[
P_\delta\cong\mathop{\mathrm{colim}}_{\alpha<\delta}P_\alpha
\]

at every nonzero limit ordinal \(\delta\leq\lambda\). If each successive
quotient \(Q_\alpha=P_{\alpha+1}/P_\alpha\) has projective dimension at
most \(n\), then \(\pd(P_\lambda)\leq n\).
\end{lemma}

\begin{proof}
First fix \(Y\in\mathcal A\) and assume only that
\(\Ext^1(Q_\alpha,Y)=0\) for all \(\alpha<\lambda\). Choose an injective
presentation

\[
0\longrightarrow Y\longrightarrow I\xrightarrow{p}C\longrightarrow0.
\]

The long exact Ext sequence shows that \(\Ext^1(Q,Y)=0\) precisely when
every map \(Q\to C\) lifts through \(p\). We will prove this lifting
property for \(P_\lambda\).

Given \(f:P_\lambda\to C\), denote its restriction to \(P_\alpha\) by
\(f_\alpha\). Construct compatible maps \(g_\alpha:P_\alpha\to I\)
satisfying \(p g_\alpha=f_\alpha\). There is a unique choice at stage zero.
At a successor stage, injectivity of \(I\) extends \(g_\alpha\) to a map
\(h:P_{\alpha+1}\to I\). This extension may not lift \(f_{\alpha+1}\),
but the discrepancy \(f_{\alpha+1}-ph\) vanishes on \(P_\alpha\).
Writing \(\pi_\alpha:P_{\alpha+1}\to Q_\alpha\) for the quotient map,
there is therefore a map \(b:Q_\alpha\to C\) such that

\[
f_{\alpha+1}-ph=b\pi_\alpha.
\]

By the assumed Ext vanishing, choose \(\widetilde b:Q_\alpha\to I\)
with \(p\widetilde b=b\). Then
\(g_{\alpha+1}=h+\widetilde b\pi_\alpha\) is both an extension of
\(g_\alpha\) and a lift of \(f_{\alpha+1}\). At a limit stage, the
compatible earlier lifts induce a map from the colimit. Their identities
with the restrictions of \(f\) persist by the colimit's universal property.
This constructs \(g_\lambda\), proving \(\Ext^1(P_\lambda,Y)=0\).

Now fix any \(d>n\) and any coefficient object \(Y\). Choose successive
injective presentations

\[
0\longrightarrow Y_j\longrightarrow I_j\longrightarrow Y_{j+1}
\longrightarrow0,
\qquad Y_0=Y,
\qquad 0\leq j<d-1.
\]

Dimension shifting gives, for every object \(Q\),

\[
\Ext^d(Q,Y)\cong\Ext^1(Q,Y_{d-1}).
\]

The same coefficient object \(Y_{d-1}\) works for all successive
quotients. Their dimension bounds make the right-hand side vanish, so the
degree-one argument applies. Shifting back gives
\(\Ext^d(P_\lambda,Y)=0\). Since \(d>n\) and \(Y\) were arbitrary,
the required projective-dimension bound follows.\leanmark{transfinite}
\end{proof}

\begin{remark}
The proof constructs compatible lifts; it does not commute Ext with an
arbitrary filtered colimit. In particular, the closure statement requires
no separate assumption that arbitrary filtered colimits are exact.
\end{remark}

\subsection{Counting compact opens}

\begin{proposition}\label{prop:compact-open-dimension}
Let \(X\) be Hausdorff and totally disconnected. If an open subset
\(W\subseteq X\) is a union of at most \(\aleph_n\) compact-open subsets
of \(X\), then \(\pd(P(W))\leq n\).
\end{proposition}

\begin{proof}
We induct on \(n\). For \(n=0\), enumerate the given compact opens as
\(K_0,K_1,\ldots\), allowing empty terms, and set

\[
L_j=K_j\setminus\bigcup_{i<j}K_i.
\]

The earlier union is finite and compact, hence closed in the Hausdorff
space \(X\); consequently \(L_j\) is open. It is also closed in the
compact set \(K_j\), since the earlier union is open. Thus the \(L_j\)
are compact open, pairwise disjoint, and have union \(W\).
Lemma~\ref{lem:projective-opens} makes \(P(W)\) projective. This includes
finite and empty covers.
Concretely, sections through a sheaf epimorphism lift separately on every
\(L_j\), and the lifts glue without compatibility conditions because the
pieces are disjoint. This explains why no finiteness condition on the
disjoint family is needed. Local compactness of the ambient space is not
required for this step or for the proposition.

Assume the assertion for \(n\). Pad a family of at most
\(\aleph_{n+1}\) compact opens with empty sets to index it by the initial
ordinal \(\kappa=\omega_{n+1}\), and put

\[
W_\alpha=\bigcup_{\beta<\alpha}K_\beta
\quad(\alpha\leq\kappa),
\qquad W_\kappa=W.
\]

These opens give a continuous filtration of free sheaves. Indeed,
inclusions induce monomorphisms: on stalks the free-open sheaf is \(\Z\)
inside its representing open and zero outside, and the nonzero inclusion
maps are identities. Moreover, the free sheaf of a directed union is the
colimit of the free sheaves of its members. By the representing property,
this assertion says exactly that compatible sections on the members glue
uniquely to a section on their union. Hence \(P(W_0)=0\), continuity holds
at limit ordinals, and the final object is \(P(W)\).

For two opens \(V,T\), the free-open Mayer--Vietoris sequence is

\[
0\longrightarrow P(V\cap T)
\longrightarrow P(V)\oplus P(T)
\longrightarrow P(V\cup T)\longrightarrow0.
\]

The first map is the signed pair of inclusions, and the second adds the
inclusions into the union. Exactness can be checked on stalks; on the
intersection the maps are \(u\mapsto(u,-u)\) and \((v,t)\mapsto v+t\).
It follows, by taking the quotient by \(P(V)\), that the successive
quotient \(Q_\alpha=P(W_{\alpha+1})/P(W_\alpha)\) fits into

\[
0\longrightarrow P(W_\alpha\cap K_\alpha)
\longrightarrow P(K_\alpha)
\longrightarrow Q_\alpha\longrightarrow0.
\]

For every \(\alpha<\kappa\), the initial-ordinal property gives
\(|\alpha|<\aleph_{n+1}\), hence \(|\alpha|\leq\aleph_n\). Since

\[
W_\alpha\cap K_\alpha
=\bigcup_{\beta<\alpha}(K_\beta\cap K_\alpha),
\]

the induction hypothesis bounds the left-hand term's projective dimension
by \(n\). The middle term is projective by
Lemma~\ref{lem:projective-opens}. For \(d>n+1\), the long exact Ext
sequence identifies \(\Ext^d(Q_\alpha,M)\) with
\(\Ext^{d-1}(P(W_\alpha\cap K_\alpha),M)\), which vanishes.
Thus \(\pd(Q_\alpha)\leq n+1\).
Lemma~\ref{lem:transfinite-dimension} completes the induction.
Only the size of each proper initial segment of \(\omega_{n+1}\) enters
the argument. There is no appeal to the continuum hypothesis or to any
extra cardinal-arithmetic assumption.\leanmark{cover-bound}
\end{proof}

Taking \(W=X\), the identification \(P(X)\cong\underline{\Z}_X\)
immediately gives the following consequence.

\begin{corollary}\label{cor:cover-vanishing}
If a Hausdorff totally disconnected space \(X\) has a compact-open cover
of cardinality at most \(\aleph_n\), then
\(H^d(X;\mathcal M)=0\) for every abelian sheaf \(\mathcal M\) and
every \(d>n\).
\end{corollary}

\subsection{Cardinality and weight}

\begin{corollary}\label{cor:universal-minimality}
Let \(V\) be locally compact, Hausdorff, and totally disconnected, let
\(r\geq1\), and let \(\mathcal M\) be any abelian sheaf on \(V\).
Then

\[
H^r(V;\mathcal M)\neq0
\quad\Longrightarrow\quad
|V|\geq\aleph_r
\quad\text{and}\quad
w(V)\geq\aleph_r.
\]

\end{corollary}

\begin{proof}
Such a space has a basis of compact opens. To recall why, given a point
and an open neighborhood, choose a compact neighborhood inside it. This
compact Hausdorff totally disconnected subspace is zero-dimensional, so
it contains a clopen neighborhood of the point lying inside its own
interior in \(V\). That smaller neighborhood is both open in \(V\) and
compact. Thus compact opens refine every neighborhood.

Choosing one compact-open
neighborhood of each point gives a cover with at most \(|V|\) members.
Thus \(|V|\leq\aleph_n\) implies vanishing above degree \(n\) by
Corollary~\ref{cor:cover-vanishing}.

For weight, choose a basis \(\mathcal B\) of cardinality \(w(V)\).
Every compact open is a finite union of basis elements: cover it by basis
neighborhoods contained in it and use compactness. Choosing one such
finite description for each compact open gives

\[
|\operatorname{CompactOpens}(V)|
\leq\max(\aleph_0,w(V)).
\]

Indeed, identical finite descriptions have identical unions, so this
choice is injective into the set of finite subsets of \(\mathcal B\).
For an infinite basis, that set has the same cardinality as the basis.
All compact opens together cover \(V\). Therefore \(w(V)\leq\aleph_n\)
also forces vanishing above degree \(n\); the displayed maximum handles
finite or empty bases as well.

Finally, a cardinal strictly below \(\aleph_r\), for positive finite
\(r\), is at most \(\aleph_{r-1}\). Either \(|V|<\aleph_r\) or
\(w(V)<\aleph_r\) would therefore imply \(H^r(V;\mathcal M)=0\).
Taking contrapositives proves both inequalities.\leanmark{minimality}
\end{proof}

A nonzero product of \(r\) degree-one classes is a nonzero degree-\(r\)
cohomology class, so these inequalities apply to every space considered in
Aoki's question. The examples of Theorem~\ref{thm:main} attain both bounds
simultaneously. This proves the asserted optimality for cardinality and
weight.
